\documentclass[10pt,a4paper]{article}
\usepackage[margin=18mm]{geometry}
\usepackage[T2A]{fontenc}
\usepackage[utf8]{inputenc}
\usepackage[english,russian]{babel}
\usepackage{amsmath,amssymb,booktabs,tabularx,graphicx,array}
\usepackage[unicode,colorlinks=true,urlcolor=blue,linkcolor=blue]{hyperref}
\setlength{\parindent}{0pt}
\setlength{\parskip}{5pt}
\emergencystretch=3em
\newcommand{\R}{\operatorname{R}}
\newcommand{\MAE}{\operatorname{MAE}}
\title{t-Flow и PFGM++: интеграция и проверка на Exp / Spiral}
\author{}
\date{11 сентября 2026}
\begin{document}
\maketitle
\vspace{-12mm}

\section*{Что проверяем и почему}
Добавлены два запрошенных метода с негауссовским распределением зашумлённого
наблюдения при фиксированной чистой точке. Вопрос: дают ли они разумное качество
LID при общих архитектуре, бюджете и выборе масштаба? Предыдущий короткий
Student-t пилот показал, что низкий MAE после выбора масштаба совместим с плохой
производной сети при малом шуме. Поэтому проверены и тестовая ошибка LID,
и расхождение с точным posterior, и denoising на свежем шуме.

\textbf{Постановка.} Использованы исходные Exp и Spiral, коэффициенты размерности30
и их канонический PCA-рендер размерности784. Истинный LID равен2 для Exp и1 для
Spiral. По уже принятому правилу оба представления этих PCA-наборов получают
общий spectral residual backbone: ширина512, четыре блока, ковариационная
оболочка по fit-данным. Её ранг3/2 означает ранг линейной оболочки Exp/Spiral,
а не подставленный истинный LID. Новые методы имеют ровно столько же параметров,
сколько Gaussian-контроль на соответствующих данных. U-Net-маршрут тоже
проверен обучением, загрузкой и производными; четыре исследуемые PCA-ячейки
не являются проверкой качества U-Net на Arrows.

На каждую из четырёх ячеек заново обучены t-Flow, PFGM++ и Gaussian posterior
rectified flow: всего12 независимых обучений, один seed на модель/ячейку.
Одинаковы99000 fit и1000 source-train holdout объектов, fit-only центрирование и
RMS-нормировка,128000 обновлений, batch256, AdamW2e-4, clipping1, последний
8000-шаговый cosine/EMA участок. Веса выбирает нативный holdout loss, а не LID.
Все1000 test объектов открываются для основного измерения после фиксации весов
и масштаба. Равны обновления и число сетевых примеров; равенство FLOPs не заявлено.

\section*{Что именно добавлено}
Оба метода используют радиальное многомерное Student-t ядро. Один chi-square
знаменатель делится всеми координатами объекта; независимые одномерные Student-t
по координатам дали бы другое распределение. У t-Flow фиксировано $\nu=5$,
у PFGM++ дополнительная размерность128, то есть эквивалентное Student-t $\nu=128$.
Это два нативных метода одной семьи ядер. Параметры хвостов одинаковы на всех
датасетах и не подбираются по LID.

Нативный t-Flow использует равномерное время и noise MSE; PFGM++ --- EDM weighted
denoising loss и lognormal noise scale. Среднее log-sigma в PFGM++ переведено
из исходного sigma-data0.5 в общие единицы sigma-data1: $-1.2+\log2$, std1.2.
Распределение условно ограничено поддержкой через inverse CDF; хвосты самого
шума не обрезаются. Это явные нативные различия методов. Такое сравнение
не изолирует только форму ядра при одинаковом loss и sampling.

Во всех таблицах $\lambda$ --- RMS шума \emph{на одну ambient-координату}
в fit-нормированных единицах. Для Student-t его нативный scale равен
$\sigma=\lambda\sqrt{(\nu-2)/\nu}$; радиус PFGM++ равен
$r=\lambda\sqrt{\nu-2}$. Старые Gaussian $\lambda$ не меняются.

Student-t шум сохраняет связь между касательными и нормальными координатами.
При данных в линейной оболочке ранга$k$, ambient-размерности$N$ и наблюдаемой
нормальной компоненте$n$ условное активное ядро имеет
\[
\nu_{\rm active}=\nu+N-k,\qquad
\lambda_{\rm active}^2=
\frac{(\nu-2)\lambda^2+\|n\|^2}{\nu+N-k-2}.
\]
Именно этот scale подаётся общему core. Так нормальная информация учитывается
без дополнительных обучаемых признаков; image U-Net получает все координаты.

Пусть$b(y)$ --- предсказанный posterior mean чистой точки. Оценивается скаляр
$\R(y,\lambda)=\operatorname{tr}D_y b(y)$, то есть trace его Jacobian.
Для нативной скорости t-Flow это$Nt+t(1-t)\operatorname{div}v$, для радиального
поля PFGM++ --- $N-r\operatorname{div}u$. Gaussian добавка квадрата posterior
displacement здесь не применяется: запросы \texttt{full}/\texttt{fm\_to\_score}
для новых методов явно отклоняются.

\section*{Выбор масштаба и тестовая ошибка}
На1000 source-train holdout объектах выбирается один масштаб, минимизирующий
$\operatorname{mean}_i|\R_i-d_i|$, по тем же29 полушагам log2 от1/256 до64.
Он фиксируется перед test. Основной pointwise Kneedle использует ровно ту же
сетку без LID-меток, с прежними параметрами и ambient fallback. Граничные
выборы не скрываются, плохие автоматические оценки не заменяются supervised.

Для прямого сравнения response дополнительно применено заранее объявленное
правило выбора к Gaussian-$\R$. Это отдельная диагностика: основной Gaussian
benchmark по-прежнему выбирает масштаб по full и хранит dependent response
на этом масштабе. Дополнительное правило использует только train-holdout,
фиксируется перед отдельным test-$\R$ и не меняет основную запись benchmark.

Ниже MAE считается по всем1000 test объектам; меньше лучше. Интервалы95\% ---
query bootstrap при фиксированных весах и выборе масштаба, не неопределённость
от повторного обучения. До запуска practical gate был задан как MAE$<0.5$
на каждой задаче/представлении; он относится к supervised finite-scale оценке.

\input{results/score_table.tex}
\input{results/outcome.tex}

\section*{Проверка того, что выучила модель}
На первых32 фиксированных source-train holdout точках вычислен posterior
\emph{непрерывных} исходных законов. Для Exp интегрируются равномерные осевая
координата и угол, с радиусом$3\exp(-u)$; для Spiral --- исходная
$\big(\sin(t^2)/t,\cos(t^2)/t\big)$ при равномерном$t\in[1,100]$.
Это не KDE по обучающим точкам. Учитывается точный измеренный Gram рендера.
Student posterior получен Gamma-смесью непрерывных Gaussian posterior,
включая производную весов смеси. Эталон не участвует в обучении.

На минимальном масштабе точный response Spiral уже может превышать1: близкие
витки не разрешаются конечным шумом. На коэффициентах средние точные значения
равны1.5043 для t-Flow и1.6178 для PFGM++; на PCA-рендере ---1.00002 и1.2497.
Из-за нормального условного распределения Student смена ambient-размерности
меняет posterior даже при почти изометрическом рендере. Поэтому расхождение
с1 и расхождение \emph{с точным posterior} --- разные измерения.

В таблице ниже $\R_{\min}$ --- средний response на32 holdout точках при
$\lambda=1/256$. Последний столбец --- среднее$|\R_{\rm model}-\R_{\rm exact}|$
на \emph{выбранном} масштабе, по тем же32 точкам. Низкий selected MAE сам по себе
не доказывает точность при$\lambda\to0$.
\input{results/oracle_table.tex}

Дополнительно128 test точек независимо зашумлены восемь раз каждая, своим
нативным шумом каждой модели. Измерен$\mathbb{E}\|b(X+\lambda Z)-X\|^2/\lambda^2$.
Сопоставлена обученная сеть с её начальным нулевым residual-head, то есть
ковариационным preconditioner без выученной нелинейной поправки. Внутри строки
используются те же чистые точки и шум; меньше лучше. Между разными ядрами
наблюдения и posterior различаются, поэтому сравнение величин риска между
методами не интерпретируется как причинный эффект одного ядра.
\input{results/denoising_table.tex}

\section*{Корректность и воспроизводимость}
Проверены радиальный закон шума, нативный loss, точное условное ядро,
field/readout identities, finite differences, фактическое обучение и загрузка
vector/U-Net, сохранённые selection receipts и полные pointwise результаты.
252 regression проверки и две дополнительные U-Net training/reload проверки
прошли.23 независимых проверок эталона дали максимальный разрыв$6.01\cdot10^{-10}$.
Во всех1392 проверках уточнения quadrature максимальное изменение response
было$6.48\cdot10^{-7}$. Повторный расчёт на фактических float32 входах меняет
эталон не более чем на$1.66\cdot10^{-5}$ на проверенных4 точках каждого
представления и29 масштабах.

Полные численные значения, исходники проверок, bootstrap выбора масштаба и
ссылки на чекпойнты находятся рядом: \texttt{results/metrics.csv},
\texttt{results/verification.json}, \texttt{protocol.md}.
Новые модели входят в общую конфигурационную матрицу39$\times$13=507;
эта работа измеряет только указанные12 обучений, не всю матрицу.

Первоисточники:
\href{https://arxiv.org/html/2410.14171v2}{Heavy-Tailed Diffusion Models,
Appendix B, Eq.164--166};
\href{https://proceedings.mlr.press/v202/xu23m/xu23m.pdf}{PFGM++, Eq.4--6};
\href{https://github.com/Newbeeer/pfgmpp/blob/main/training/loss.py}{официальный EDM loss PFGM++}.

\clearpage
\begin{center}
\includegraphics[width=\linewidth]{results/mean_response_logx.pdf}
\end{center}
Средние по одним32 holdout точкам: сплошные линии --- обученные модели,
пунктир --- соответствующий точный posterior. Точки на линиях показывают
масштаб, выбранный по всем1000 train-holdout объектам. Логарифмическая шкала
показывает также малые шумы.

\clearpage
\begin{center}
\includegraphics[width=\linewidth]{results/mean_response_linearx.pdf}
\end{center}
Те же значения на линейной шкале$\lambda$. Ни точки, ни кандидаты выбора
масштаба не изменены.

\clearpage
\begin{center}
\includegraphics[width=\linewidth,height=.89\textheight,keepaspectratio]{results/point_response_linearx.pdf}
\end{center}
Одна фиксированная holdout точка в каждой ячейке. Её ID указан на панели,
и все модели внутри ячейки используют одну и ту же точку.
\end{document}
